\chapter{De onde vêm as ideias da JVM}
\label{cap:historia-vms}

\proximasaulas

Bytecode de pilha, JIT, inline cache, desotimização, metadados ricos no
arquivo de classe: quase nada disso nasceu com o Java. Este capítulo
visita cinco sistemas, em ordem aproximada, e liga cada um a algo que a
Parte~\ref{parte:java-jvm} mostrou na HotSpot de hoje. A ideia não é fazer arqueologia: é
entender por que a JVM é como é, e o que ela poderia ter sido.

\begin{description}
\item[Pascal-P e o p-code] (anos~1970) Compilar para uma máquina que não
  existe e portar só o interpretador.
\item[Smalltalk-80] (1983) Uma VM especificada em livro, com bytecodes,
  interpretador e memória de objetos.
\item[Self] (1989--1994) O laboratório do JIT adaptativo, dos inline
  caches polimórficos e da desotimização.
\item[CLR] (anos~2000) A VM da Microsoft, que fez outra escolha para
  generics.
\item[Dalvik e ART] (Android) Uma VM de registradores e, depois, a mistura
  de AOT com JIT.
\end{description}

Não incluímos um sistema quando não conseguimos uma fonte primária para
ele; por isso a lista é curta.

\section{Pascal-P e o p-code}
\label{sec:hist-pcode}

\paragraph{A ideia.}
O livro de \textcite{livro:pemberton1982}, que comenta linha a linha o
compilador Pascal-P4, resume o desenho logo na introdução: em vez de
gerar código para uma máquina real, o compilador gera o que veio a ser
chamado de \termo{p-code}{\emph{Pseudo-code}: instruções de uma máquina
hipotética, a P-machine, executadas por um interpretador.}, \enquote{for a
hypothetical stack-based computer that is in many ways ideal for Pascal
compilation}. O compilador e o interpretador do p-code são escritos em
Pascal.

\paragraph{Por que isso portava.}
Para levar o sistema a uma máquina nova, o livro descreve caminhos como
este: compilar o compilador Pascal-P com qualquer compilador Pascal já
existente, usá-lo para traduzir o próprio compilador para p-code e, daí
em diante, só traduzir à mão o \emph{interpretador} para a máquina nova.
O livro diz que o Pascal-P foi a base de muitos sistemas Pascal,
\enquote{including the well-known UCSD system}.

\paragraph{A máquina.}
O cap.~10 descreve a P-machine: um processador e uma memória, em que
muitas operações usam a pilha, chamadas de procedimento empilham
parâmetros e endereço de retorno, e um conjunto pequeno de registradores
(contador de programa, ponteiro de pilha e outros) controla a execução.

\paragraph{O que isso explica na JVM.}
\begin{itemize}
\item \emph{Bytecode de pilha.} A JVM também é uma máquina de pilha:
  \texttt{iload}, \texttt{iadd}, \texttt{ireturn} empilham e desempilham
  operandos, como a seção~\ref{sec:javap-basico} mostrou no
  \texttt{javap}. A troca, segundo \textcite{paper:vm-showdown}: código
  de pilha é menor, mas precisa de mais instruções para a mesma conta
  (seção~\ref{sec:hist-dalvik}).
\item \emph{Portar a VM, não o programa.} O \enquote{escreva uma vez, rode
  em qualquer lugar} do Java é a mesma aposta: o \texttt{.class} é o
  mesmo, e o trabalho de cada plataforma é ter uma JVM.
\item \emph{O custo.} Interpretar tem um preço por instrução:
  \textcite{paper:vm-showdown} descrevem cada despacho de instrução como
  um desvio indireto caro e, em geral, imprevisível. A JVM paga isso só
  no começo e depois compila; as seções~\ref{sec:hist-smalltalk}
  e~\ref{sec:hist-self} mostram de onde veio essa saída.
\end{itemize}

\section{Smalltalk-80: a VM especificada em livro}
\label{sec:hist-smalltalk}

\paragraph{Imagem e máquina virtual.}
O \enquote{Blue Book}, \citetitle{livro:goldberg1983}
\parencite{livro:goldberg1983}, dedica a quarta parte à implementação. O
capítulo~26 separa dois componentes: a \emph{virtual image}, com todos
os objetos do sistema, e a \emph{virtual machine}, que dá dinâmica a
eles. A tarefa de quem implementa o Smalltalk-80 numa máquina nova é
escrever a máquina virtual; a imagem é carregada nela.

\paragraph{Bytecodes de pilha.}
O mesmo capítulo descreve o caminho do código: o compilador traduz cada
método para \enquote{sequences of eight-bit instructions called
bytecodes}, que são instruções para um interpretador orientado a pilha.
O exemplo do livro é o método \texttt{center} de \texttt{Rectangle}
(\texttt{\^{}origin + corner / 2}), que vira seis bytecodes: empilha a
primeira variável de instância, empilha a segunda, envia \texttt{+},
empilha a constante 2, envia \texttt{/} e retorna o topo da pilha. Abaixo
do interpretador fica a \emph{object memory}, que guarda os objetos; os
capítulos~27 a~30 especificam interpretador e memória em detalhe.

Compare com qualquer listagem \texttt{javap} da Parte~II: instruções de um
byte, operandos na pilha, chamadas de método como instrução
(\texttt{invokevirtual}) e um heap com coleta de lixo por baixo. A JVM
não copiou o Smalltalk, mas a forma é a mesma.

\paragraph{Inline caches.}
O Smalltalk também foi onde a chamada de método ficou rápida.
\textcite{paper:self-pics} contam a ideia, atribuída a
\textcite{paper:smalltalk-80-deutsch-schiffman}: o tipo do receptor num
ponto de chamada quase nunca muda, então o endereço do método encontrado
pode ser guardado no próprio ponto de chamada, e o método chamado só
confere se o tipo ainda é o mesmo. Esse
\termo{inline cache}{Cache do resultado da busca de método guardado no
próprio ponto de chamada, no código do chamador.} acerta 95\% das vezes em
código Smalltalk, número que o paper reúne de três trabalhos anteriores
(um deles o de Deutsch e Schiffman).

\paragraph{O que isso explica na JVM.}
\begin{itemize}
\item O formato da VM: bytecode compacto de pilha, interpretador e
  memória de objetos gerenciada. A seção~\ref{sec:javap-basico} é a
  versão Java do capítulo~26 do Blue Book.
\item As chamadas virtuais baratas. Quando a seção~\ref{sec:jit} fala de
  pontos de chamada monomórficos e megamórficos, está falando de inline
  caches e do que acontece quando eles não bastam, o problema que o Self
  atacou em seguida.
\end{itemize}

\section{Self e a origem do JIT da HotSpot}
\label{sec:hist-self}

O Self era uma linguagem dinâmica, sem classes, em que tudo é envio de
mensagem. Isso deveria torná-la lenta. O grupo de Stanford (e depois da
Sun) passou anos provando que não precisava ser, e as técnicas que
inventou para isso estão na HotSpot.

\paragraph{1989: customização e compilação dinâmica.}
\textcite{paper:self-efficient-impl} agrupam objetos clonados do mesmo
protótipo em \emph{maps} (o que dá ao compilador informação de tipo) e
compilam em runtime várias versões de um mesmo método, cada uma
especializada para o \emph{map} do receptor, o que permite fazer inline
das mensagens enviadas a \texttt{self}. Os autores afirmam que essa
implementação roda duas vezes mais rápido que o Smalltalk mais rápido da
época.

\paragraph{1991: inline caches polimórficos.}
O inline cache da seção anterior falha quando um ponto de chamada vê
vários tipos. \textcite{paper:self-pics} propõem os
\termo{polymorphic inline caches}{PICs: inline caches com uma entrada
por tipo de receptor já visto naquele ponto de chamada.}, e notam o
efeito colateral mais importante: o PIC \emph{registra} os tipos que
apareceram em cada ponto de chamada, e o compilador pode usar isso ao
recompilar. Números dos autores, medidos no próprio sistema Self: mediana
de 11\% de ganho só com PICs e de 27\% quando o compilador usa os tipos
coletados, com cerca de metade das chamadas não inlinadas eliminadas.

\paragraph{1992: desotimização.}
Código muito otimizado é difícil de depurar.
\textcite{paper:self-deoptimization} resolvem isso
\enquote{dynamically deoptimizing code on demand}: só as ativações que
estão sendo depuradas voltam a código não otimizado, e o resto roda em
velocidade máxima. Com isso, também dá para mudar um programa em
execução e ver o efeito na hora. Uma ferramenta de depuração virou a
\termo{deoptimization}{Desotimização: trocar, em plena execução, frames de
código otimizado por frames equivalentes de código não otimizado.} que a
HotSpot usa quando uma especulação falha
(\citetitle{artigo:shipilev-uncommon-traps}).

\paragraph{1994: otimização adaptativa.}
A tese de \textcite{paper:holzle-thesis-adaptive-optimization} junta as
peças: um compilador rápido gera o código inicial, as partes muito usadas
são recompiladas por um compilador otimizador usando o \emph{type
feedback} e o código não otimizado é recriado quando necessário. Leia de
novo a seção~\ref{sec:jit}: interpretador ou C1 primeiro, C2 com perfil
depois, \emph{uncommon trap} e desotimização quando o perfil mente. É o
mesmo desenho.

\paragraph{A ligação, nas palavras de quem fez.}
Não precisamos adivinhar a ligação. O paper do compilador C2,
\textcite{paper:hotspot-server-compiler}, diz que o runtime usa
otimização adaptativa para concentrar a compilação nos métodos críticos
e cita Hölzle e Ungar para isso; e diz que, quando o carregamento de uma
classe invalida uma decisão de inlining, os métodos dependentes são
desotimizados, citando o paper de 1992 acima.
\textcite{paper:truffle-partial-evaluation}, da Oracle Labs, abrem o
paper dizendo que VMs como a HotSpot e o V8 \enquote{follow the design
that was first implemented for the SELF language}.

\paragraph{O que isso explica na JVM.}
\begin{itemize}
\item Por que o warm-up existe e por que a HotSpot coleta perfil: o
  compilador otimizador só é bom porque viu os tipos reais
  (seção~\ref{sec:jit}).
\item Por que chamadas virtuais e interfaces não são caras na prática:
  inline caches e inlining guiado por tipo, herdados do Self.
\item Por que carregar uma classe nova pode deixar o programa mais lento
  por um instante: desotimização de métodos que assumiam um só
  implementador (seção~\ref{sec:class-loading}).
\item Por que o nativo pode perder em throughput de pico, como na
  medição do Quarkus (seção~\ref{sec:jvm-vs-nativo}), embora o paper do
  Native Image fale em preservar o desempenho de pico: sem perfil da
  execução real e sem desotimização, o compilador só pode apostar no que
  prova no build (seção~\ref{sec:vms-jit-aot}).
\end{itemize}

\section{CLR: a VM que fez generics de outro jeito}
\label{sec:hist-clr}

\paragraph{O padrão.}
A plataforma .NET é especificada pelo padrão
\citetitle{doc:ecma335} (6ª edição, junho de 2012, equivalente ao
ISO/IEC~23271). A página da Ecma descreve a
\termo{CLI}{\emph{Common Language Infrastructure}: o padrão que define
como programas em várias linguagens rodam sobre uma mesma VM.} como uma
infraestrutura para que programas em várias linguagens rodem em vários
ambientes sem serem reescritos, e lista as partes do padrão: a
\emph{Virtual Execution System} (VES, a máquina virtual), os metadados
e o seu formato de arquivo, e o conjunto de instruções
\emph{Common Intermediate Language} (CIL). O CLR é a implementação da
Microsoft. Troque os nomes e a JVM aparece: a especificação da JVM
(\citetitle{doc:jvms25}) define a máquina, o formato \texttt{.class} e o
conjunto de instruções.

\paragraph{A diferença que importa: generics.}
Em Java, generics são resolvidos pelo compilador e apagados: a JLS chama
isso de \emph{type erasure}, \enquote{a mapping from types (possibly
including parameterized types and type variables) to types (that are
never parameterized types or type variables)}
(\citetitle{doc:jls25}, §4.6). Em runtime, \texttt{List<String>} e
\texttt{List<Integer>} são a mesma classe.

O CLR fez outra escolha. \textcite{paper:clr-generics} descrevem generics
com \enquote{exact run-time types}: a VM sabe, em runtime, que uma lista
é de inteiros. A implementação aproveita o JIT para fazer
\enquote{just-in-time type specialization}, com código compartilhado
entre instanciações de mesma representação; os primeiros resultados
sugerem que quem usa generics não paga sobrecusto por isso.

\paragraph{O que isso explica na JVM.}
\begin{itemize}
\item Por que não existe \texttt{new T[]} nem \texttt{instanceof
  List<String>} em Java: a informação foi apagada antes de chegar à VM.
\item Por que \texttt{List<Integer>} guarda objetos \texttt{Integer} e
  não \texttt{int}: com erasure, a VM só vê \texttt{Object}, e o
  autoboxing da seção~\ref{sec:autoboxing} é o preço. Uma VM que conhece
  o tipo exato pode especializar o código para o tipo primitivo, como o
  paper do CLR descreve.
\item Que erasure é uma escolha da plataforma Java, não um limite de
  toda VM com bytecode: o CLR mostra a alternativa, com o custo de levar
  os tipos genéricos para dentro da VM e do JIT.
\end{itemize}

\section{Dalvik e ART: registradores, AOT e JIT}
\label{sec:hist-dalvik}

\paragraph{Dalvik: uma VM de registradores para Java.}
O Android usa Java como linguagem, mas não roda uma JVM: o runtime
executa bytecode no formato DEX (\citetitle{doc:android-art}).
A especificação diz: \enquote{The machine is
register-based, and frames are fixed in size upon creation}; cada frame
tem um número fixo de registradores de 32~bits, valores de 64~bits usam
pares, e o desenho tenta \enquote{approximately imitate common real
architectures and C-style calling conventions}
(\citetitle{doc:android-dalvik-bytecode}).

\paragraph{Pilha ou registradores?}
\textcite{paper:vm-showdown} mediram a troca que o Dalvik fez. Eles
traduziram bytecode JVM de pilha para um formato de registradores dentro
de uma JVM interpretada (Sun J2ME CDC~1.0), com SPECjvm98 e Java Grande.
Resultado, medido pelos autores (Trinity College Dublin e TU Wien, 2005):
a versão de registradores executa, em média, mais de 47\% menos
instruções da VM, com bytecode cerca de 25\% maior; num Pentium~4, leva
em média 32,3\% menos tempo com despacho por \texttt{switch} em C e cerca
de 26,5\% menos com \emph{threaded dispatch}. Note o que foi medido: um
\emph{interpretador}. Num JIT, o bytecode é só a entrada do compilador,
que gera código para os registradores da máquina real; o C2, por
exemplo, faz alocação de registradores por coloração de grafos
\parencite{paper:hotspot-server-compiler}. A vantagem medida vale para a
fase interpretada, não para o código quente.

\paragraph{ART: AOT na instalação, JIT com perfil.}
O sucessor do Dalvik é o ART. A documentação diz que \enquote{ART and
Dalvik are compatible runtimes running DEX bytecode} e que, na
instalação, o ART compila os apps com a ferramenta \texttt{dex2oat}; o GC
é descrito como majoritariamente concorrente, com uma só pausa
(\citetitle{doc:android-art}). O ART também tem um JIT
\enquote{with code profiling}: um método pode estar interpretado,
compilado pelo JIT ou compilado AOT, e um daemon usa o perfil coletado
pelo JIT para guiar a compilação AOT, o que a documentação chama de
\emph{profile-guided compilation} (\citetitle{doc:android-art-jit}).

\paragraph{O que isso explica na JVM.}
\begin{itemize}
\item Que o bytecode de pilha da JVM é um formato de distribuição, não
  de execução: Dalvik e ART pegam o mesmo código Java e o executam em
  outra forma.
\item Que AOT e JIT não são rivais. O ART usa o perfil do JIT para
  compilar antes; no OpenJDK, o Projeto Leyden vai na mesma direção com o
  cache AOT da seção~\ref{sec:jvm-vs-nativo} e o perfil de métodos
  antecipado do \citetitle{doc:jep515}.
\end{itemize}

\section{Onde aprofundar}

\subsection*{Livros}
\citetitle{livro:pemberton1982} \parencite{livro:pemberton1982}, com o
texto completo publicado pelo autor, é a melhor leitura sobre o p-code:
compilador e interpretador comentados. \citetitle{livro:goldberg1983}
\parencite{livro:goldberg1983}, capítulos~26 a~30, especifica a VM do
Smalltalk-80.

\subsection*{Papers}
\small
\begin{longtable}{@{}p{5.2cm} p{7cm}@{}}
\toprule
Fonte & Para quê \\
\midrule \endhead
\textcite{paper:aycock-history-jit} & Survey de quarenta anos de JIT,
  de 1960 a 2000: o mapa para o resto da lista \\
\textcite{paper:self-pics} & Inline caches polimórficos; curto e com PDF
  aberto \\
\textcite{paper:self-deoptimization} & De onde vem a desotimização \\
\textcite{paper:holzle-thesis-adaptive-optimization} & Otimização
  adaptativa completa: o desenho da HotSpot, em 1994 \\
\textcite{paper:hotspot-server-compiler} & O C2 descrito por quem o
  fez; PDF aberto na USENIX \\
\textcite{paper:clr-generics} & Generics com tipos exatos em runtime, o
  contraponto ao erasure \\
\textcite{paper:vm-showdown} & Pilha contra registradores, com dados \\
\bottomrule
\end{longtable}
\normalsize

\subsection*{Documentação}
\citetitle{doc:ecma335} (o padrão da CLI, com download do PDF na página
da Ecma), \citetitle{doc:android-dalvik-bytecode},
\citetitle{doc:android-art} e \citetitle{doc:android-art-jit}.

\subsection*{De quem fez o Java}
Nomes, bios e mais obras de cada pessoa no apêndice~\ref{ap:pessoas}.
\begin{itemize}
\obra*{paper:java-language-environment}{O white paper de 1996: o que o
  Java herdou e o que recusou das linguagens e VMs deste capítulo.}
\obra*{paper:feel-of-java}{Gosling sobre as decisões de design, em
  primeira pessoa.}
\obra*{livro:lindholm2014}{O item~1.1, \enquote{A Bit of History}, conta
  de onde veio a JVM.}
\obra*{video:java-story}{O documentário oficial, com quem construiu o
  Java.}
\end{itemize}

\subsection*{Pesquisa acadêmica}
Itens que completam os papers acima. Acesso e links no
apêndice~\ref{ap:pesquisa}.
\begin{itemize}
\pesquisa{paper:pascal-p-code-generation}{O compilador Pascal-P de Zurique
  e a máquina de pilha hipotética que a UCSD adaptou no p-System.}
\pesquisa{paper:self-power-of-simplicity}{A linguagem Self: protótipos,
  slots e mensagens, o começo da linha Self, Strongtalk, HotSpot.}
\pesquisa{paper:self-type-feedback}{Type feedback, a ideia central dos JITs
  modernos.}
\pesquisa{paper:jalapeno-adaptive-optimization}{A otimização adaptativa da
  JVM de pesquisa da IBM, que virou a Jikes RVM.}
\pesquisa{paper:lua50-implementation}{A primeira VM de registradores de uso
  amplo, feita na PUC-Rio.}
\pesquisa*{paper:hazan-maquina-registradores-java}{Pilha ou registradores
  para Java, com o olhar de dispositivos móveis (PUC-Rio).}
\pesquisa{paper:jit-value-specialization}{Pesquisa da UFMG em JIT
  especulativo: especializar funções JavaScript para os valores reais.}
\pesquisa{paper:isamap-dbt}{Tradução dinâmica de binários do IC-Unicamp,
  parente do JIT: traduzir e guardar código em execução.}
\end{itemize}

\section{Desafios práticos}

\desafio{Pegue o método \texttt{center} do Smalltalk-80 da
seção~\ref{sec:hist-smalltalk}, escreva o equivalente em Java
(\texttt{record Retangulo(Ponto origem, Ponto canto)}) e compare a saída
do \texttt{javap -c} com os seis bytecodes do Blue Book. O que a JVM
precisa dizer explicitamente que o Smalltalk não precisa?}

\desafio{Escreva um ponto de chamada que vê primeiro um só tipo e,
depois de aquecido, passa a ver três. Rode com
\texttt{-XX:+PrintCompilation} e procure a linha \enquote{made not
entrant}. Relacione o que você viu com os PICs e a desotimização do
Self.}

\desafio{Escreva um método genérico que tenta criar um \texttt{T[]} e
veja o erro do compilador. Leia a §4.6 da JLS e explique por que o
mesmo código seria possível com os generics do CLR.}

\desafio{Leia o cap.~10 de Pemberton e Daniels e escreva, em
Java, um interpretador de umas dez instruções de uma máquina de pilha.
Meça quanto tempo ele leva para somar um milhão de números e compare
com o mesmo laço em Java puro. Onde está a diferença?}

